\chapter{Dynamic Programming: DP on Strings}
\chapterepigraph{Two strings rarely need a third dimension. Almost
everything in this chapter is answered by a single 2D table, indexed by
one pointer into each string, where the recurrence is just: do the
characters at these two pointers match, and if so, how does that
agreement combine with the answer for both pointers advanced together?}

\section{One Table, Two Pointers, Two Strings}
\label{sec:dp-strings-intro}

Every problem in this chapter compares two strings $s_1$ (length $n$)
and $s_2$ (length $m$) using a DP state $(i, j)$, where $i$ indexes a
prefix (or suffix) of $s_1$ and $j$ indexes a prefix (or suffix) of
$s_2$. The recurrence almost always branches on a single question: does
$s_1[i] = s_2[j]$?

\begin{patternbox}[The Core Recurrence Almost Every Problem Specializes]
\begin{itemize}
  \item \textbf{If $s_1[i] = s_2[j]$:} this matched pair of characters
        is worth keeping, so add $1$ (or otherwise credit the match) and
        recurse on \emph{both} pointers advanced:
        $f(i,j) = 1 + f(i{-}1,j{-}1)$ (for subsequence-length problems),
        or a similarly structured combination for other problems.
  \item \textbf{If $s_1[i] \ne s_2[j]$:} the characters cannot both be
        used in this match, so the best answer comes from advancing
        \emph{one} pointer at a time and taking whichever is better:
        $f(i,j) = \max(f(i{-}1,j),\ f(i,j{-}1))$, or, for edit-distance-style
        problems, the minimum over several specific single-step or
        dual-step operations (insert, delete, replace).
\end{itemize}
The entire chapter is variations on this one template: \textbf{Longest
Common Subsequence} establishes it in its cleanest form; every other
section either reuses it directly (Shortest Common Supersequence,
Distinct Subsequences), reduces to it (Longest Palindromic Subsequence,
Delete Operations, Minimum Insertions for Palindrome), or extends it with
a small number of additional operations (Edit Distance, Wildcard
Matching).
\end{patternbox}

\subsection*{Why Recursing on a Suffix-or-Prefix Pair, Not a Single Index, Is the New Idea}

Every problem in the DP 1D and DP Grids chapters so far has used at most
two coupled positions (a row and column on one shared structure, or two
agents on the same row). Here, the two indices $i$ and $j$ belong to
\textbf{independent} strings -- $s_1[i]$ and $s_2[j]$ are not from the
same array at all, so there is no inherent geometric relationship between
them the way there was for grid DP's row/column. The table is still
filled exactly like a grid (row by row, or with both loops nested), but
the meaning of ``the cell at $(i,j)$'' is now ``comparing a prefix of
$s_1$ against a prefix of $s_2$,'' a genuinely new state shape worth
internalizing carefully in the chapter's first section before relying on
it everywhere else.

\newpage

\section{Longest Common Subsequence}
\label{sec:dp-lcs}

\problemstatement{Given two strings $s_1$ (length $n$) and $s_2$ (length
$m$), find the length of their \textbf{longest common subsequence}
(LCS): the longest sequence of characters that appears, in order but not
necessarily contiguously, in both strings.}

\begin{patternbox}
This is the chapter's foundational recurrence, established in its
cleanest form: \emph{``longest sequence appearing in order in both
strings''} is answered by comparing characters one pair at a time from
the end (or start) of each string, branching on whether they match.
\end{patternbox}

\subsection*{Deriving the recurrence}

It is conventional in string DP to use a \textbf{1-indexed} table, where
$f(i,j)$ denotes the answer for the prefixes $s_1[0..i-1]$ (length $i$)
and $s_2[0..j-1]$ (length $j$) -- this lets $f(0,j)$ and $f(i,0)$ cleanly
represent ``one of the strings is empty,'' whose LCS length is trivially
$0$, without needing special-case index arithmetic for negative indices.

\[
  f(i,j) =
  \begin{cases}
    0 & \text{if } i=0 \text{ or } j=0, \\
    1 + f(i{-}1,j{-}1) & \text{if } s_1[i{-}1] = s_2[j{-}1], \\
    \max(f(i{-}1,j),\ f(i,j{-}1)) & \text{otherwise.}
  \end{cases}
\]

If the last characters of both prefixes match, that character is part of
some longest common subsequence, so we credit it and recurse on both
prefixes shortened by one. If they don't match, the last character of
\emph{at least one} prefix is not part of the LCS (it cannot be part of
\emph{both} prefixes' contribution simultaneously, since they disagree),
so we try discarding each one in turn and take the better result.

\subsection*{Approach 1: Recursion}

\begin{lstlisting}
#include <bits/stdc++.h>
using namespace std;

int lcsRecursive(int i, int j, string& s1, string& s2) {
    if (i == 0 || j == 0) return 0;

    if (s1[i - 1] == s2[j - 1]) {
        return 1 + lcsRecursive(i - 1, j - 1, s1, s2);
    }
    return max(lcsRecursive(i - 1, j, s1, s2), lcsRecursive(i, j - 1, s1, s2));
}

int main() {
    string s1 = "abcde", s2 = "ace";
    cout << lcsRecursive(s1.size(), s2.size(), s1, s2) << endl; // 3
    return 0;
}
\end{lstlisting}

\begin{complexitybox}[Approach 1 Complexity]
\textbf{Time.} $O(2^{n+m})$ -- two branches per mismatch, depth up to
$n+m$.
\textbf{Space.} $O(n+m)$ recursion stack.
\end{complexitybox}

\subsection*{Approach 2: Memoization}

\begin{lstlisting}
#include <bits/stdc++.h>
using namespace std;

int lcsMemo(int i, int j, string& s1, string& s2, vector<vector<int>>& memo) {
    if (i == 0 || j == 0) return 0;
    if (memo[i][j] != -1) return memo[i][j];

    if (s1[i - 1] == s2[j - 1]) {
        return memo[i][j] = 1 + lcsMemo(i - 1, j - 1, s1, s2, memo);
    }
    return memo[i][j] = max(lcsMemo(i - 1, j, s1, s2, memo),
                             lcsMemo(i, j - 1, s1, s2, memo));
}

int main() {
    string s1 = "abcde", s2 = "ace";
    int n = s1.size(), m = s2.size();
    vector<vector<int>> memo(n + 1, vector<int>(m + 1, -1));
    cout << lcsMemo(n, m, s1, s2, memo) << endl; // 3
    return 0;
}
\end{lstlisting}

\begin{complexitybox}[Approach 2 Complexity]
\textbf{Time.} $O(n \cdot m)$.
\textbf{Space.} $O(n \cdot m)$ memo $+$ $O(n+m)$ recursion stack.
\end{complexitybox}

\subsection*{Approach 3: Tabulation}

\begin{lstlisting}
#include <bits/stdc++.h>
using namespace std;

int lcsTabulation(string& s1, string& s2) {
    int n = s1.size(), m = s2.size();
    vector<vector<int>> dp(n + 1, vector<int>(m + 1, 0));

    for (int i = 1; i <= n; i++) {
        for (int j = 1; j <= m; j++) {
            if (s1[i - 1] == s2[j - 1]) {
                dp[i][j] = 1 + dp[i - 1][j - 1];
            } else {
                dp[i][j] = max(dp[i - 1][j], dp[i][j - 1]);
            }
        }
    }
    return dp[n][m];
}

int main() {
    string s1 = "abcde", s2 = "ace";
    cout << lcsTabulation(s1, s2) << endl; // 3
    return 0;
}
\end{lstlisting}

\subsection*{Dry run of the tabulation table}

\dryrun{Take $s_1 = \texttt{"abcde"}$, $s_2 = \texttt{"ace"}$.}

\begin{center}
\begin{tabular}{c|c|c|c|c|c|c}
 & \texttt{""} & \texttt{a} & \texttt{c} & \texttt{e} \\ \hline
\texttt{""} & 0 & 0 & 0 & 0 \\ \hline
\texttt{a} & 0 & 1 & 1 & 1 \\ \hline
\texttt{b} & 0 & 1 & 1 & 1 \\ \hline
\texttt{c} & 0 & 1 & 2 & 2 \\ \hline
\texttt{d} & 0 & 1 & 2 & 2 \\ \hline
\texttt{e} & 0 & 1 & 2 & 3
\end{tabular}
\end{center}

At $(i,j)$ corresponding to \texttt{c} (row) and \texttt{c} (column):
match, so $\textit{dp}=1+\textit{dp}[\text{prev row}][\text{prev col}]=1+1=2$.
At the final cell (\texttt{e}, \texttt{e}): match, $\textit{dp}=1+2=3$.
The LCS length is $3$ (the subsequence \texttt{"ace"} itself).

\begin{complexitybox}[Approach 3 Complexity]
\textbf{Time.} $O(n \cdot m)$.
\textbf{Space.} $O(n \cdot m)$ table, no recursion stack.
\end{complexitybox}

\subsection*{Approach 4: Space Optimization}

\begin{lstlisting}
#include <bits/stdc++.h>
using namespace std;

int lcsSpaceOptimized(string& s1, string& s2) {
    int n = s1.size(), m = s2.size();
    vector<int> prev(m + 1, 0), current(m + 1, 0);

    for (int i = 1; i <= n; i++) {
        for (int j = 1; j <= m; j++) {
            if (s1[i - 1] == s2[j - 1]) {
                current[j] = 1 + prev[j - 1];
            } else {
                current[j] = max(prev[j], current[j - 1]);
            }
        }
        prev = current;
    }
    return prev[m];
}

int main() {
    string s1 = "abcde", s2 = "ace";
    cout << lcsSpaceOptimized(s1, s2) << endl; // 3
    return 0;
}
\end{lstlisting}

\begin{complexitybox}[Approach 4 Complexity]
\textbf{Time.} $O(n \cdot m)$.
\textbf{Space.} $O(m)$ for two rolling rows.
\end{complexitybox}

\begin{takeawaybox}
This match/mismatch recurrence -- credit a match and advance both
pointers, or try discarding one character from each side on a mismatch
-- is the single template the rest of this chapter builds from. Memorize
its shape (not just its code) before moving on, since every later
section either reuses it verbatim on a transformed input (reversing one
string, concatenating two strings) or extends its mismatch case with
extra operations.
\end{takeawaybox}

\section{Print the Longest Common Subsequence}
\label{sec:dp-print-lcs}

\problemstatement{Given two strings $s_1$ and $s_2$, construct (not just
measure the length of) one valid longest common subsequence.}

\begin{patternbox}
\emph{``Print/construct the actual solution, not just its length/count''}
is a new keyword pattern this chapter introduces: once a tabulation table
is built, \textbf{walking it backward} from the final answer cell,
re-deriving at each step which recurrence case produced that cell's
value, reconstructs an actual optimal solution -- without needing any new
DP derivation, only a careful reverse walk through the already-computed
table.
\end{patternbox}

\subsection*{Understanding the reconstruction}

The Section~\ref{sec:dp-lcs} tabulation table $\textit{dp}[i][j]$ must be
built in full first (this requires the complete $O(nm)$ table, not the
space-optimized $O(m)$ rolling-row version, since reconstruction needs
to look back arbitrarily far). Starting at $\textit{dp}[n][m]$, walk
backward: if $s_1[i-1] = s_2[j-1]$, this character is part of the LCS --
record it and move diagonally to $(i-1,j-1)$. Otherwise, move to
whichever of $(i-1,j)$ or $(i,j-1)$ has the larger $\textit{dp}$ value
(matching whichever branch the forward recurrence actually took to
produce $\textit{dp}[i][j]$'s value). Continue until reaching a row or
column of $0$. The characters recorded, read in reverse order (since we
walked backward from the end), form one valid LCS.

\subsection*{Why retracing the recurrence's decision backward is correct}

The tabulation table already encodes, implicitly, which choice (match,
or which discard) was responsible for each cell's value -- we simply
need to recompute that same comparison once more during the backward
walk, since the table only stores the resulting \emph{length}, not which
branch produced it. Because the forward recurrence's correctness was
already established in Section~\ref{sec:dp-lcs}, redoing its exact
decision rule in reverse is guaranteed to retrace a genuine optimal path
through the table, character by character, until the base case (an
empty prefix) is reached.

\subsection*{Algorithm}

\begin{enumerate}
  \item Build the full $\textit{dp}$ table as in
        Section~\ref{sec:dp-lcs}'s tabulation.
  \item $i \gets n$, $j \gets m$, $\textit{result} \gets \texttt{""}$.
  \item While $i>0$ and $j>0$:
  \begin{enumerate}
    \item If $s_1[i-1] = s_2[j-1]$: prepend $s_1[i-1]$ to
          $\textit{result}$; $i \gets i-1$, $j \gets j-1$.
    \item Else if $\textit{dp}[i-1][j] \ge \textit{dp}[i][j-1]$:
          $i \gets i-1$.
    \item Else: $j \gets j-1$.
  \end{enumerate}
  \item Return \textit{result}.
\end{enumerate}

\subsection*{Dry run}

\dryrun{Using Section~\ref{sec:dp-lcs}'s table for $s_1=\texttt{"abcde"}$,
$s_2=\texttt{"ace"}$.}

\begin{itemize}
  \item Start at $(i,j)=(5,3)$: $s_1[4]=\texttt{'e'}$,
        $s_2[2]=\texttt{'e'}$, match. Record \texttt{'e'}. Move to
        $(4,2)$.
  \item $(4,2)$: $s_1[3]=\texttt{'d'}$, $s_2[1]=\texttt{'c'}$, no
        match. $\textit{dp}[3][2]=2 \ge \textit{dp}[4][1]=1$: move to
        $(3,2)$.
  \item $(3,2)$: $s_1[2]=\texttt{'c'}$, $s_2[1]=\texttt{'c'}$, match.
        Record \texttt{'c'}. Move to $(2,1)$.
  \item $(2,1)$: $s_1[1]=\texttt{'b'}$, $s_2[0]=\texttt{'a'}$, no
        match. $\textit{dp}[1][1]=1 \ge \textit{dp}[2][0]=0$: move to
        $(1,1)$.
  \item $(1,1)$: $s_1[0]=\texttt{'a'}$, $s_2[0]=\texttt{'a'}$, match.
        Record \texttt{'a'}. Move to $(0,0)$. Loop ends ($i=0$).
\end{itemize}

Characters recorded, in the order found (backward): \texttt{e, c, a}.
Reversing gives \texttt{"ace"} -- the LCS itself.

\subsection*{C++ implementation}

\begin{lstlisting}
#include <bits/stdc++.h>
using namespace std;

string printLCS(string& s1, string& s2) {
    int n = s1.size(), m = s2.size();
    vector<vector<int>> dp(n + 1, vector<int>(m + 1, 0));

    for (int i = 1; i <= n; i++) {
        for (int j = 1; j <= m; j++) {
            if (s1[i - 1] == s2[j - 1]) {
                dp[i][j] = 1 + dp[i - 1][j - 1];
            } else {
                dp[i][j] = max(dp[i - 1][j], dp[i][j - 1]);
            }
        }
    }

    string result;
    int i = n, j = m;
    while (i > 0 && j > 0) {
        if (s1[i - 1] == s2[j - 1]) {
            result.push_back(s1[i - 1]);
            i--; j--;
        } else if (dp[i - 1][j] >= dp[i][j - 1]) {
            i--;
        } else {
            j--;
        }
    }
    reverse(result.begin(), result.end());
    return result;
}

int main() {
    string s1 = "abcde", s2 = "ace";
    cout << printLCS(s1, s2) << endl; // ace
    return 0;
}
\end{lstlisting}

\begin{complexitybox}
\textbf{Time Complexity.} $O(n \cdot m)$ to build the table, plus
$O(n+m)$ for the backward walk -- dominated by the table construction.
\textbf{Space Complexity.} $O(n \cdot m)$ for the full table --
\textbf{not} reducible to a rolling row, since reconstruction needs
access to arbitrary earlier rows during the backward walk.
\end{complexitybox}

\begin{takeawaybox}
``Print/construct the solution'' problems always need the \emph{full}
tabulation table kept in memory, sacrificing the space optimization that
a length/count-only version of the same problem could use -- the
backward walk that reconstructs an answer needs to revisit cells the
rolling-row trick would have already discarded. This trade-off (full
table for reconstruction vs.\ rolling row for the value alone) is worth
remembering whenever a problem explicitly asks for the solution itself.
\end{takeawaybox}

\section{Longest Common Substring}
\label{sec:dp-longest-common-substring}

\problemstatement{Given two strings $s_1$ and $s_2$, find the length of
their longest common \textbf{substring}: a longest run of
\textbf{contiguous} characters that appears, unbroken, in both strings.}

\begin{patternbox}
The keyword shift from \emph{subsequence} (Section~\ref{sec:dp-lcs}) to
\emph{substring} -- contiguity now required -- breaks the mismatch case
of the LCS recurrence entirely: a substring cannot ``skip'' a
mismatched character and continue, so a mismatch must \textbf{reset} the
current run to length $0$, rather than falling back to the best of two
neighboring sub-answers.
\end{patternbox}

\subsection*{Deriving the recurrence}

Let $f(i,j)$ be the length of the common substring \textbf{ending
exactly at} $s_1[i-1]$ and $s_2[j-1]$ (not the best answer over all
prefixes, as in LCS -- specifically the run ending right here):
\[
  f(i,j) =
  \begin{cases}
    1 + f(i{-}1,j{-}1) & \text{if } s_1[i{-}1] = s_2[j{-}1], \\
    0 & \text{otherwise.}
  \end{cases}
\]
Because the answer is not necessarily at $f(n,m)$ (the longest common
run could end anywhere in either string), we additionally track a
running global maximum over every $f(i,j)$ computed.

\subsection*{Approach 1: Recursion}

\begin{ideabox}[Why a Direct Recursive Translation Is Awkward Here]
Unlike every previous problem in this book, $f(i,j)$ alone does not
directly answer the original question -- the answer is
$\max_{i,j} f(i,j)$, a quantity that is naturally accumulated during a
forward (tabulation-style) sweep, not easily expressed as a single
top-down recursive call. For this reason, Longest Common Substring is
typically presented starting from tabulation, with the recursive/memoized
forms understood as computing the same per-cell $f(i,j)$ values, just
collected via top-down calls instead of a bottom-up loop.
\end{ideabox}

\subsection*{Approach 2: Memoization}

\begin{lstlisting}
#include <bits/stdc++.h>
using namespace std;

int f(int i, int j, string& s1, string& s2, vector<vector<int>>& memo, int& best) {
    if (i == 0 || j == 0) return 0;
    if (memo[i][j] != -1) return memo[i][j];

    int result = 0;
    if (s1[i - 1] == s2[j - 1]) {
        result = 1 + f(i - 1, j - 1, s1, s2, memo, best);
    }
    best = max(best, result);
    return memo[i][j] = result;
}

int longestCommonSubstringMemo(string& s1, string& s2) {
    int n = s1.size(), m = s2.size();
    vector<vector<int>> memo(n + 1, vector<int>(m + 1, -1));
    int best = 0;

    for (int i = 1; i <= n; i++) {
        for (int j = 1; j <= m; j++) {
            f(i, j, s1, s2, memo, best);
        }
    }
    return best;
}

int main() {
    string s1 = "abcde", s2 = "abfce";
    cout << longestCommonSubstringMemo(s1, s2) << endl; // 2
    return 0;
}
\end{lstlisting}

\begin{complexitybox}[Approach 2 Complexity]
\textbf{Time.} $O(n \cdot m)$.
\textbf{Space.} $O(n \cdot m)$ memo.
\end{complexitybox}

\subsection*{Approach 3: Tabulation}

\begin{lstlisting}
#include <bits/stdc++.h>
using namespace std;

int longestCommonSubstringTabulation(string& s1, string& s2) {
    int n = s1.size(), m = s2.size();
    vector<vector<int>> dp(n + 1, vector<int>(m + 1, 0));
    int best = 0;

    for (int i = 1; i <= n; i++) {
        for (int j = 1; j <= m; j++) {
            if (s1[i - 1] == s2[j - 1]) {
                dp[i][j] = 1 + dp[i - 1][j - 1];
            } else {
                dp[i][j] = 0;
            }
            best = max(best, dp[i][j]);
        }
    }
    return best;
}

int main() {
    string s1 = "abcde", s2 = "abfce";
    cout << longestCommonSubstringTabulation(s1, s2) << endl; // 2
    return 0;
}
\end{lstlisting}

\subsection*{Dry run of the tabulation table}

\dryrun{Take $s_1=\texttt{"abcde"}$, $s_2=\texttt{"abfce"}$.}

\begin{center}
\begin{tabular}{c|c|c|c|c|c|c}
 & \texttt{""} & \texttt{a} & \texttt{b} & \texttt{f} & \texttt{c} & \texttt{e} \\ \hline
\texttt{""} & 0 & 0 & 0 & 0 & 0 & 0 \\ \hline
\texttt{a} & 0 & 1 & 0 & 0 & 0 & 0 \\ \hline
\texttt{b} & 0 & 0 & 2 & 0 & 0 & 0 \\ \hline
\texttt{c} & 0 & 0 & 0 & 0 & 1 & 0 \\ \hline
\texttt{d} & 0 & 0 & 0 & 0 & 0 & 0 \\ \hline
\texttt{e} & 0 & 0 & 0 & 0 & 0 & 1
\end{tabular}
\end{center}

The running maximum reaches $2$ at cell (\texttt{b}, \texttt{b}) --
the common substring \texttt{"ab"} -- and never exceeds it elsewhere
(every other match is isolated, contributing only length $1$ runs).

\begin{complexitybox}[Approach 3 Complexity]
\textbf{Time.} $O(n \cdot m)$.
\textbf{Space.} $O(n \cdot m)$ table, no recursion stack.
\end{complexitybox}

\subsection*{Approach 4: Space Optimization}

\begin{lstlisting}
#include <bits/stdc++.h>
using namespace std;

int longestCommonSubstringSpaceOptimized(string& s1, string& s2) {
    int n = s1.size(), m = s2.size();
    vector<int> prev(m + 1, 0), current(m + 1, 0);
    int best = 0;

    for (int i = 1; i <= n; i++) {
        for (int j = 1; j <= m; j++) {
            if (s1[i - 1] == s2[j - 1]) {
                current[j] = 1 + prev[j - 1];
            } else {
                current[j] = 0;
            }
            best = max(best, current[j]);
        }
        prev = current;
    }
    return best;
}

int main() {
    string s1 = "abcde", s2 = "abfce";
    cout << longestCommonSubstringSpaceOptimized(s1, s2) << endl; // 2
    return 0;
}
\end{lstlisting}

\begin{complexitybox}[Approach 4 Complexity]
\textbf{Time.} $O(n \cdot m)$.
\textbf{Space.} $O(m)$ for two rolling rows.
\end{complexitybox}

\begin{takeawaybox}
``Subsequence'' (gaps allowed) and ``substring'' (no gaps allowed) are
the single most important keyword distinction in string DP, and they
produce genuinely different recurrences -- LCS falls back to a
neighboring best-of-two on a mismatch, while Longest Common Substring
resets to zero, because a broken run cannot be patched back together.
Always check this one word before assuming a string problem is ``LCS in
disguise.''
\end{takeawaybox}

\section{Shortest Common Supersequence}
\label{sec:dp-shortest-common-supersequence}

\problemstatement{Given two strings $s_1$ (length $n$) and $s_2$ (length
$m$), construct the \textbf{shortest} string that has \textbf{both}
$s_1$ and $s_2$ as subsequences.}

\begin{patternbox}
\emph{``Shortest string containing both as subsequences''} is the
``opposite'' framing of Section~\ref{sec:dp-lcs}'s LCS, and the two are
linked by a clean formula: every character the two strings \emph{share}
in their LCS needs to appear only \textbf{once} in the supersequence
(since it serves both strings' subsequence requirement simultaneously),
while every other character of each string must still appear --
collapsing the problem into one already-solved LCS computation plus a
reconstruction walk almost identical to Section~\ref{sec:dp-print-lcs}'s.
\end{patternbox}

\subsection*{Deriving the length formula}

If we simply concatenated $s_1$ and $s_2$, the result trivially contains
both as subsequences, with length $n+m$ -- but every character that is
part of \emph{some} common subsequence is being counted \textbf{twice}
in this concatenation (once from each string), when it only needs to
appear once. The longest common subsequence is exactly the largest set
of characters we can ``merge'' this way, so:
\[
  |\textit{SCS}| = n + m - |\textit{LCS}(s_1, s_2)|.
\]

\subsection*{Constructing the actual supersequence}

Run the same backward walk as Section~\ref{sec:dp-print-lcs}, but
instead of recording \emph{only} matched characters, record
\textbf{every} character visited along the way: on a match, record the
(shared) character once and move diagonally; on a mismatch, record
whichever of $s_1[i-1]$ or $s_2[j-1]$ corresponds to the direction moved
(following the same $\textit{dp}[i-1][j]$ vs.\ $\textit{dp}[i][j-1]$
comparison as before), and move only that one pointer. Once either
pointer reaches $0$, append the entirety of whatever remains of the
other string (every leftover character must still appear, in order,
since that string's subsequence requirement is otherwise unmet).

\subsection*{Algorithm}

\begin{enumerate}
  \item Build the LCS tabulation table for $s_1, s_2$ (as in
        Section~\ref{sec:dp-lcs}).
  \item $i \gets n$, $j \gets m$, $\textit{result} \gets \texttt{""}$.
  \item While $i>0$ and $j>0$:
  \begin{enumerate}
    \item If $s_1[i-1]=s_2[j-1]$: prepend $s_1[i-1]$; $i{-}{-}$,
          $j{-}{-}$.
    \item Else if $\textit{dp}[i-1][j] \ge \textit{dp}[i][j-1]$:
          prepend $s_1[i-1]$; $i{-}{-}$.
    \item Else: prepend $s_2[j-1]$; $j{-}{-}$.
  \end{enumerate}
  \item Prepend any remaining prefix of $s_1$ (if $i>0$) or $s_2$ (if
        $j>0$).
  \item Return \textit{result}.
\end{enumerate}

\subsection*{Dry run}

\dryrun{Take $s_1=\texttt{"abac"}$, $s_2=\texttt{"cab"}$.}

LCS of \texttt{"abac"} and \texttt{"cab"} is \texttt{"ab"} (length $2$),
so $|\textit{SCS}| = 4+3-2 = 5$. Constructing: walking backward through
the LCS table merges the two strings, interleaving non-matched
characters with the two shared \texttt{a} and \texttt{b}, producing one
valid result such as \texttt{"cabac"} (length $5$) -- which indeed
contains \texttt{"abac"} as a subsequence (positions $1,3,4,5$... more
directly, removing the leading \texttt{c} from \texttt{"cabac"} leaves
\texttt{"abac"} exactly) and \texttt{"cab"} as a subsequence (the first
three characters).

\subsection*{C++ implementation}

\begin{lstlisting}
#include <bits/stdc++.h>
using namespace std;

string shortestCommonSupersequence(string& s1, string& s2) {
    int n = s1.size(), m = s2.size();
    vector<vector<int>> dp(n + 1, vector<int>(m + 1, 0));

    for (int i = 1; i <= n; i++) {
        for (int j = 1; j <= m; j++) {
            if (s1[i - 1] == s2[j - 1]) {
                dp[i][j] = 1 + dp[i - 1][j - 1];
            } else {
                dp[i][j] = max(dp[i - 1][j], dp[i][j - 1]);
            }
        }
    }

    string result;
    int i = n, j = m;
    while (i > 0 && j > 0) {
        if (s1[i - 1] == s2[j - 1]) {
            result.push_back(s1[i - 1]);
            i--; j--;
        } else if (dp[i - 1][j] >= dp[i][j - 1]) {
            result.push_back(s1[i - 1]);
            i--;
        } else {
            result.push_back(s2[j - 1]);
            j--;
        }
    }
    while (i > 0) { result.push_back(s1[i - 1]); i--; }
    while (j > 0) { result.push_back(s2[j - 1]); j--; }

    reverse(result.begin(), result.end());
    return result;
}

int main() {
    string s1 = "abac", s2 = "cab";
    string scs = shortestCommonSupersequence(s1, s2);
    cout << scs << " (length " << scs.size() << ")" << endl;
    return 0;
}
\end{lstlisting}

\begin{complexitybox}
\textbf{Time Complexity.} $O(n \cdot m)$ to build the LCS table, plus
$O(n+m)$ for the reconstruction walk.
\textbf{Space Complexity.} $O(n \cdot m)$ for the full table (needed for
reconstruction, as in Section~\ref{sec:dp-print-lcs}).
\end{complexitybox}

\begin{takeawaybox}
Shortest Common Supersequence is not a new DP recurrence at all -- it is
LCS's length formula, $n+m-|\textit{LCS}|$, plus a reconstruction walk
that is a small variation on Print LCS's (recording every visited
character, not only the matches). Whenever a problem can be phrased as
``the complement or merge of an already-solved quantity,'' check the
arithmetic relationship before reaching for a new table.
\end{takeawaybox}

\section{Distinct Subsequences}
\label{sec:dp-distinct-subsequences}

\problemstatement{Given strings $s$ (length $n$) and $t$ (length $m$),
count the number of distinct ways $t$ can be formed as a subsequence of
$s$ (counting different choices of \emph{positions} in $s$ as distinct,
even if they spell the same characters).}

\begin{patternbox}
\emph{``Count the number of ways one string appears as a subsequence of
another''} reuses Section~\ref{sec:dp-lcs}'s match/mismatch branching,
but on a \textbf{match}, there are now \textbf{two} valid choices to
\emph{sum} rather than one branch to take unconditionally: use the
current character of $s$ to match $t$, \emph{or} skip it and hope a
later character of $s$ matches instead -- both contribute to the count,
since $s$ may contain the needed character multiple times.
\end{patternbox}

\subsection*{Deriving the recurrence}

Let $f(i,j)$ be the number of ways $t[0..j-1]$ (length $j$) can be
formed as a subsequence of $s[0..i-1]$ (length $i$).

\[
  f(i,j) =
  \begin{cases}
    f(i{-}1,j{-}1) + f(i{-}1,j) & \text{if } s[i{-}1]=t[j{-}1], \\
    f(i{-}1,j) & \text{otherwise.}
  \end{cases}
\]

If the characters match, we may either \textbf{use} this occurrence in
$s$ to match $t$'s current character (advancing both, $f(i{-}1,j{-}1)$),
\emph{or} \textbf{skip} it, leaving $t$'s current character to be
matched by some earlier-yet-unconsidered part of $s$ ($f(i{-}1,j)$) --
both possibilities lead to valid, distinct subsequence selections, so
their counts are summed. If the characters don't match, this position
in $s$ cannot contribute to matching $t[j-1]$ at all, so only the skip
option remains.

Base cases: $f(i,0)=1$ for every $i$ (the empty string is a subsequence
of anything, in exactly one way -- selecting nothing). $f(0,j)=0$ for
every $j>0$ (a non-empty $t$ cannot be a subsequence of an empty $s$).

\subsection*{Approach 1: Recursion}

\begin{lstlisting}
#include <bits/stdc++.h>
using namespace std;

int distinctSubsequencesRecursive(int i, int j, string& s, string& t) {
    if (j == 0) return 1;
    if (i == 0) return 0;

    if (s[i - 1] == t[j - 1]) {
        return distinctSubsequencesRecursive(i - 1, j - 1, s, t) +
               distinctSubsequencesRecursive(i - 1, j, s, t);
    }
    return distinctSubsequencesRecursive(i - 1, j, s, t);
}

int main() {
    string s = "babgbag", t = "bag";
    cout << distinctSubsequencesRecursive(s.size(), t.size(), s, t) << endl; // 5
    return 0;
}
\end{lstlisting}

\begin{complexitybox}[Approach 1 Complexity]
\textbf{Time.} $O(2^n)$ in the worst case.
\textbf{Space.} $O(n)$ recursion stack.
\end{complexitybox}

\subsection*{Approach 2: Memoization}

\begin{lstlisting}
#include <bits/stdc++.h>
using namespace std;

int distinctSubsequencesMemo(int i, int j, string& s, string& t,
                              vector<vector<long long>>& memo) {
    if (j == 0) return 1;
    if (i == 0) return 0;
    if (memo[i][j] != -1) return memo[i][j];

    long long result;
    if (s[i - 1] == t[j - 1]) {
        result = distinctSubsequencesMemo(i - 1, j - 1, s, t, memo) +
                 distinctSubsequencesMemo(i - 1, j, s, t, memo);
    } else {
        result = distinctSubsequencesMemo(i - 1, j, s, t, memo);
    }
    return memo[i][j] = result;
}

int main() {
    string s = "babgbag", t = "bag";
    int n = s.size(), m = t.size();
    vector<vector<long long>> memo(n + 1, vector<long long>(m + 1, -1));
    cout << distinctSubsequencesMemo(n, m, s, t, memo) << endl; // 5
    return 0;
}
\end{lstlisting}

\begin{complexitybox}[Approach 2 Complexity]
\textbf{Time.} $O(n \cdot m)$.
\textbf{Space.} $O(n \cdot m)$ memo $+$ $O(n)$ recursion stack.
\end{complexitybox}

\subsection*{Approach 3: Tabulation}

\begin{lstlisting}
#include <bits/stdc++.h>
using namespace std;

long long distinctSubsequencesTabulation(string& s, string& t) {
    int n = s.size(), m = t.size();
    vector<vector<long long>> dp(n + 1, vector<long long>(m + 1, 0));

    for (int i = 0; i <= n; i++) dp[i][0] = 1;

    for (int i = 1; i <= n; i++) {
        for (int j = 1; j <= m; j++) {
            if (s[i - 1] == t[j - 1]) {
                dp[i][j] = dp[i - 1][j - 1] + dp[i - 1][j];
            } else {
                dp[i][j] = dp[i - 1][j];
            }
        }
    }
    return dp[n][m];
}

int main() {
    string s = "babgbag", t = "bag";
    cout << distinctSubsequencesTabulation(s, t) << endl; // 5
    return 0;
}
\end{lstlisting}

\subsection*{Dry run of the tabulation table}

\dryrun{Take $s=\texttt{"babgbag"}$, $t=\texttt{"bag"}$.}

The five ways \texttt{"bag"} appears as a subsequence of
\texttt{"babgbag"} correspond to choosing one of the two \texttt{'b'}s
(positions $0,2,4$... specifically $3$ valid \texttt{'b'} positions
combined with valid later \texttt{'a'},\texttt{'g'} pairs) in
combination with the available \texttt{'a'} and \texttt{'g'}
occurrences after it; building the table row by row accumulates these
combinations automatically via the sum rule, reaching
$\textit{dp}[7][3]=5$.

\begin{complexitybox}[Approach 3 Complexity]
\textbf{Time.} $O(n \cdot m)$.
\textbf{Space.} $O(n \cdot m)$ table, no recursion stack.
\end{complexitybox}

\subsection*{Approach 4: Space Optimization}

\begin{lstlisting}
#include <bits/stdc++.h>
using namespace std;

long long distinctSubsequencesSpaceOptimized(string& s, string& t) {
    int n = s.size(), m = t.size();
    vector<long long> prev(m + 1, 0), current(m + 1, 0);
    prev[0] = current[0] = 1;

    for (int i = 1; i <= n; i++) {
        for (int j = 1; j <= m; j++) {
            if (s[i - 1] == t[j - 1]) {
                current[j] = prev[j - 1] + prev[j];
            } else {
                current[j] = prev[j];
            }
        }
        prev = current;
    }
    return prev[m];
}

int main() {
    string s = "babgbag", t = "bag";
    cout << distinctSubsequencesSpaceOptimized(s, t) << endl; // 5
    return 0;
}
\end{lstlisting}

\begin{complexitybox}[Approach 4 Complexity]
\textbf{Time.} $O(n \cdot m)$.
\textbf{Space.} $O(m)$ for two rolling rows.
\end{complexitybox}

\begin{takeawaybox}
A character match does not always mean ``take it unconditionally and
advance both pointers'' -- when the same character could plausibly be
matched by \emph{several} different positions in the longer string (as
here, where $s$ may contain repeats), a match should branch into
\textbf{both} ``use it'' and ``save it for later, skip for now,''
summing their counts. This is a subtler version of the match case than
Section~\ref{sec:dp-lcs}'s, worth contrasting carefully.
\end{takeawaybox}

\section{Edit Distance}
\label{sec:dp-edit-distance}

\problemstatement{Given two strings $s_1$ (length $n$) and $s_2$ (length
$m$), find the minimum number of operations -- \textbf{insert},
\textbf{delete}, or \textbf{replace} a single character -- needed to
transform $s_1$ into $s_2$.}

\begin{patternbox}
This is Section~\ref{sec:dp-lcs}'s match/mismatch template, extended:
the \textbf{match} case is unchanged (no operation needed, advance both
pointers for free), but the \textbf{mismatch} case is no longer a simple
``discard one side'' -- it now branches into \textbf{three} possible
operations, each costing $1$ and corresponding to a different way of
advancing the two pointers.
\end{patternbox}

\subsection*{Deriving the recurrence}

\[
  f(i,j) =
  \begin{cases}
    f(i{-}1,j{-}1) & \text{if } s_1[i{-}1]=s_2[j{-}1], \\
    1 + \min\big(f(i{-}1,j{-}1),\ f(i{-}1,j),\ f(i,j{-}1)\big) & \text{otherwise.}
  \end{cases}
\]

On a mismatch, three operations are available, each mapping to a
specific pointer movement: \textbf{replace} $s_1[i-1]$ with
$s_2[j-1]$ (both characters consumed, advance both pointers --
$f(i{-}1,j{-}1)$); \textbf{delete} $s_1[i-1]$ (advance only $s_1$'s
pointer, leaving $s_2[j-1]$ still to be matched -- $f(i{-}1,j)$); or
\textbf{insert} a copy of $s_2[j-1]$ into $s_1$ (this newly-inserted
character matches $s_2[j-1]$ for free, so only $s_2$'s pointer advances
-- $f(i,j{-}1)$). We pay $1$ for whichever operation is used, plus the
cost of solving the resulting smaller sub-problem, and take the cheapest
option.

Base cases: $f(0,j)=j$ (transform an empty $s_1$ into $s_2$'s first $j$
characters by $j$ insertions), $f(i,0)=i$ (transform $s_1$'s first $i$
characters into an empty string by $i$ deletions).

\subsection*{Approach 1: Recursion}

\begin{lstlisting}
#include <bits/stdc++.h>
using namespace std;

int editDistanceRecursive(int i, int j, string& s1, string& s2) {
    if (i == 0) return j;
    if (j == 0) return i;

    if (s1[i - 1] == s2[j - 1]) {
        return editDistanceRecursive(i - 1, j - 1, s1, s2);
    }
    int replace = editDistanceRecursive(i - 1, j - 1, s1, s2);
    int del     = editDistanceRecursive(i - 1, j, s1, s2);
    int insert  = editDistanceRecursive(i, j - 1, s1, s2);
    return 1 + min({replace, del, insert});
}

int main() {
    string s1 = "horse", s2 = "ros";
    cout << editDistanceRecursive(s1.size(), s2.size(), s1, s2) << endl; // 3
    return 0;
}
\end{lstlisting}

\begin{complexitybox}[Approach 1 Complexity]
\textbf{Time.} $O(3^{n+m})$ in the worst case -- three branches per
mismatch.
\textbf{Space.} $O(n+m)$ recursion stack.
\end{complexitybox}

\subsection*{Approach 2: Memoization}

\begin{lstlisting}
#include <bits/stdc++.h>
using namespace std;

int editDistanceMemo(int i, int j, string& s1, string& s2,
                      vector<vector<int>>& memo) {
    if (i == 0) return j;
    if (j == 0) return i;
    if (memo[i][j] != -1) return memo[i][j];

    if (s1[i - 1] == s2[j - 1]) {
        return memo[i][j] = editDistanceMemo(i - 1, j - 1, s1, s2, memo);
    }
    int replace = editDistanceMemo(i - 1, j - 1, s1, s2, memo);
    int del     = editDistanceMemo(i - 1, j, s1, s2, memo);
    int insert  = editDistanceMemo(i, j - 1, s1, s2, memo);
    return memo[i][j] = 1 + min({replace, del, insert});
}

int main() {
    string s1 = "horse", s2 = "ros";
    int n = s1.size(), m = s2.size();
    vector<vector<int>> memo(n + 1, vector<int>(m + 1, -1));
    cout << editDistanceMemo(n, m, s1, s2, memo) << endl; // 3
    return 0;
}
\end{lstlisting}

\begin{complexitybox}[Approach 2 Complexity]
\textbf{Time.} $O(n \cdot m)$.
\textbf{Space.} $O(n \cdot m)$ memo $+$ $O(n+m)$ recursion stack.
\end{complexitybox}

\subsection*{Approach 3: Tabulation}

\begin{lstlisting}
#include <bits/stdc++.h>
using namespace std;

int editDistanceTabulation(string& s1, string& s2) {
    int n = s1.size(), m = s2.size();
    vector<vector<int>> dp(n + 1, vector<int>(m + 1, 0));

    for (int j = 0; j <= m; j++) dp[0][j] = j;
    for (int i = 0; i <= n; i++) dp[i][0] = i;

    for (int i = 1; i <= n; i++) {
        for (int j = 1; j <= m; j++) {
            if (s1[i - 1] == s2[j - 1]) {
                dp[i][j] = dp[i - 1][j - 1];
            } else {
                dp[i][j] = 1 + min({dp[i - 1][j - 1], dp[i - 1][j], dp[i][j - 1]});
            }
        }
    }
    return dp[n][m];
}

int main() {
    string s1 = "horse", s2 = "ros";
    cout << editDistanceTabulation(s1, s2) << endl; // 3
    return 0;
}
\end{lstlisting}

\subsection*{Dry run of the tabulation table}

\dryrun{Take $s_1=\texttt{"horse"}$, $s_2=\texttt{"ros"}$.}

\begin{center}
\begin{tabular}{c|c|c|c|c}
 & \texttt{""} & \texttt{r} & \texttt{o} & \texttt{s} \\ \hline
\texttt{""} & 0 & 1 & 2 & 3 \\ \hline
\texttt{h} & 1 & 1 & 2 & 3 \\ \hline
\texttt{o} & 2 & 2 & 1 & 2 \\ \hline
\texttt{r} & 3 & 2 & 2 & 2 \\ \hline
\texttt{s} & 4 & 3 & 3 & 2 \\ \hline
\texttt{e} & 5 & 4 & 4 & 3
\end{tabular}
\end{center}

The answer is $\textit{dp}[5][3]=3$: one optimal sequence is
\texttt{"horse"} $\to$ \texttt{"rorse"} (replace \texttt{h}$\to$\texttt{r})
$\to$ \texttt{"rose"} (delete the second \texttt{r}) $\to$
\texttt{"ros"} (delete \texttt{e}): $3$ operations.

\begin{complexitybox}[Approach 3 Complexity]
\textbf{Time.} $O(n \cdot m)$.
\textbf{Space.} $O(n \cdot m)$ table, no recursion stack.
\end{complexitybox}

\subsection*{Approach 4: Space Optimization}

\begin{lstlisting}
#include <bits/stdc++.h>
using namespace std;

int editDistanceSpaceOptimized(string& s1, string& s2) {
    int n = s1.size(), m = s2.size();
    vector<int> prev(m + 1), current(m + 1);
    for (int j = 0; j <= m; j++) prev[j] = j;

    for (int i = 1; i <= n; i++) {
        current[0] = i;
        for (int j = 1; j <= m; j++) {
            if (s1[i - 1] == s2[j - 1]) {
                current[j] = prev[j - 1];
            } else {
                current[j] = 1 + min({prev[j - 1], prev[j], current[j - 1]});
            }
        }
        prev = current;
    }
    return prev[m];
}

int main() {
    string s1 = "horse", s2 = "ros";
    cout << editDistanceSpaceOptimized(s1, s2) << endl; // 3
    return 0;
}
\end{lstlisting}

\begin{complexitybox}[Approach 4 Complexity]
\textbf{Time.} $O(n \cdot m)$.
\textbf{Space.} $O(m)$ for two rolling rows.
\end{complexitybox}

\begin{takeawaybox}
Edit Distance generalizes the LCS template's mismatch case from ``pick
the better of two discards'' to ``pick the cheapest of three specific
operations,'' each mapping to a distinct, intuitive pointer movement.
Once this three-way branch is understood, Delete Operations
(Section~\ref{sec:dp-delete-operations}) is recoverable as the special
case where only the delete operation is permitted.
\end{takeawaybox}

\section{Delete Operations for Two Strings}
\label{sec:dp-delete-operations}

\problemstatement{Given two strings $s_1$ (length $n$) and $s_2$ (length
$m$), find the minimum number of character \textbf{deletions} (from
either string) needed to make the two strings \textbf{equal}.}

\begin{patternbox}
This is Section~\ref{sec:dp-edit-distance}'s Edit Distance restricted to
\textbf{only the delete operation} -- and that restriction collapses it
into a direct reduction to Section~\ref{sec:dp-lcs}'s LCS, rather than
needing its own three-way recurrence at all.
\end{patternbox}

\subsection*{Understanding the reduction}

If the only allowed operation is deletion, the two strings can only ever
become equal by each reducing down to some common subsequence of both --
deleting everything that isn't part of that shared subsequence. To
minimize total deletions, we want to keep as much as possible, which
means keeping the \textbf{longest} common subsequence and deleting
everything else:
\[
  \textit{minDeletions} = (n - |\textit{LCS}(s_1,s_2)|) + (m -
  |\textit{LCS}(s_1,s_2)|).
\]
The first term is how many characters of $s_1$ are deleted (everything
not part of the kept LCS); the second is the same for $s_2$.

\subsection*{All Four Approaches}

Every stage reuses Section~\ref{sec:dp-lcs}'s LCS implementation
unchanged; only this final arithmetic combination is new.

\begin{lstlisting}
#include <bits/stdc++.h>
using namespace std;

int lcsSpaceOptimized(string& s1, string& s2) {
    int n = s1.size(), m = s2.size();
    vector<int> prev(m + 1, 0), current(m + 1, 0);

    for (int i = 1; i <= n; i++) {
        for (int j = 1; j <= m; j++) {
            if (s1[i - 1] == s2[j - 1]) {
                current[j] = 1 + prev[j - 1];
            } else {
                current[j] = max(prev[j], current[j - 1]);
            }
        }
        prev = current;
    }
    return prev[m];
}

int minDeletions(string& s1, string& s2) {
    int lcsLen = lcsSpaceOptimized(s1, s2);
    return (s1.size() - lcsLen) + (s2.size() - lcsLen);
}

int main() {
    string s1 = "sea", s2 = "eat";
    cout << minDeletions(s1, s2) << endl; // 2
    return 0;
}
\end{lstlisting}

\subsection*{Dry run}

\dryrun{Take $s_1=\texttt{"sea"}$, $s_2=\texttt{"eat"}$.}

LCS of \texttt{"sea"} and \texttt{"eat"} is \texttt{"ea"} (length $2$).
$\textit{minDeletions} = (3-2)+(3-2) = 1+1 = 2$: delete \texttt{'s'}
from $s_1$ (leaving \texttt{"ea"}) and delete \texttt{'t'} from $s_2$
(leaving \texttt{"ea"}) -- both strings become \texttt{"ea"} with $2$
total deletions.

\begin{complexitybox}
\textbf{Time Complexity.} $O(n \cdot m)$, inherited directly from LCS.
\textbf{Space Complexity.} $O(m)$ for the rolling-row LCS.
\end{complexitybox}

\begin{takeawaybox}
Whenever an edit-distance-style problem restricts the allowed operations
to a strict subset (here, deletions only), check whether the restricted
version reduces to a simpler, already-solved problem rather than
re-deriving a new three-way (or however-many-way) recurrence. ``Delete
only, from both strings, to reach equality'' is exactly ``maximize the
shared kept subsequence,'' i.e.\ LCS, in disguise.
\end{takeawaybox}

\section{Longest Palindromic Subsequence}
\label{sec:dp-longest-palindromic-subsequence}

\problemstatement{Given a string $s$, find the length of its
\textbf{longest palindromic subsequence} (a subsequence that reads the
same forward and backward).}

\begin{patternbox}
\emph{``Longest subsequence that is a palindrome''} reduces directly to
Section~\ref{sec:dp-lcs}: a palindrome reads identically forward and
backward, so a subsequence of $s$ is a palindrome exactly when it is
\emph{also} a subsequence of $s$'s \textbf{reverse} -- making the longest
palindromic subsequence of $s$ exactly the longest common subsequence of
$s$ and $\textit{reverse}(s)$.
\end{patternbox}

\subsection*{Understanding the reduction}

Let $r = \textit{reverse}(s)$. Any common subsequence of $s$ and $r$
is, by construction, a sequence of characters appearing in order in $s$
that \emph{also} appears in order in $r$ -- but appearing in order in $r$
is the same as appearing in \emph{reverse} order in $s$. A sequence that
is both a subsequence of $s$ in forward order and a subsequence of $s$
in reverse order is precisely a palindromic subsequence of $s$. So:
\[
  |\textit{LPS}(s)| = |\textit{LCS}(s, \textit{reverse}(s))|.
\]

\subsection*{All Four Approaches}

Every stage is Section~\ref{sec:dp-lcs}'s LCS implementation, called
with $s_1 = s$ and $s_2 = \textit{reverse}(s)$.

\begin{lstlisting}
#include <bits/stdc++.h>
using namespace std;

int lcsSpaceOptimized(string& s1, string& s2) {
    int n = s1.size(), m = s2.size();
    vector<int> prev(m + 1, 0), current(m + 1, 0);

    for (int i = 1; i <= n; i++) {
        for (int j = 1; j <= m; j++) {
            if (s1[i - 1] == s2[j - 1]) {
                current[j] = 1 + prev[j - 1];
            } else {
                current[j] = max(prev[j], current[j - 1]);
            }
        }
        prev = current;
    }
    return prev[m];
}

int longestPalindromicSubsequence(string s) {
    string r = s;
    reverse(r.begin(), r.end());
    return lcsSpaceOptimized(s, r);
}

int main() {
    cout << longestPalindromicSubsequence("bbbab") << endl; // 4
    return 0;
}
\end{lstlisting}

\subsection*{Dry run}

\dryrun{Take $s = \texttt{"bbbab"}$.}

$r = \texttt{"babbb"}$. The LCS of \texttt{"bbbab"} and \texttt{"babbb"}
is \texttt{"bbbb"} (length $4$), which is indeed a palindrome and a
subsequence of the original \texttt{"bbbab"} (drop the \texttt{'a'}).

\begin{complexitybox}
\textbf{Time Complexity.} $O(n^2)$ (since $|s_1|=|s_2|=n$), inherited
directly from LCS.
\textbf{Space Complexity.} $O(n)$ for the rolling-row LCS.
\end{complexitybox}

\begin{takeawaybox}
``Longest palindromic X'' problems frequently reduce to ``longest common
X between the string and its reverse'' -- a reduction worth checking
first before writing a dedicated palindrome recurrence (which does also
exist for this specific problem, expanding outward from the center, but
the LCS reduction shown here requires no new derivation at all).
\end{takeawaybox}

\section{Minimum Insertions to Make a String Palindrome}
\label{sec:dp-min-insertions-palindrome}

\problemstatement{Given a string $s$ of length $n$, find the minimum
number of characters that must be \textbf{inserted} (anywhere in the
string) to make it a palindrome.}

\begin{patternbox}
A second reduction layered directly on top of
Section~\ref{sec:dp-longest-palindromic-subsequence}: the characters
\emph{not} part of the longest palindromic subsequence are exactly the
ones that need a ``partner'' inserted to balance them into a palindrome,
so the answer is simply $n$ minus the LPS length.
\end{patternbox}

\subsection*{Understanding the reduction}

Any palindromic subsequence of $s$, of length $k$, can be extended to a
full palindrome covering all of $s$ by inserting, for every character of
$s$ \emph{not} part of that subsequence, a matching mirrored character
on the opposite side. The longest palindromic subsequence minimizes the
number of characters left over needing such a partner, so:
\[
  \textit{minInsertions} = n - |\textit{LPS}(s)|.
\]

\subsection*{All Four Approaches}

Every stage reuses Section~\ref{sec:dp-longest-palindromic-subsequence}'s
implementation (itself a reduction to LCS) unchanged.

\begin{lstlisting}
#include <bits/stdc++.h>
using namespace std;

int lcsSpaceOptimized(string& s1, string& s2) {
    int n = s1.size(), m = s2.size();
    vector<int> prev(m + 1, 0), current(m + 1, 0);

    for (int i = 1; i <= n; i++) {
        for (int j = 1; j <= m; j++) {
            if (s1[i - 1] == s2[j - 1]) {
                current[j] = 1 + prev[j - 1];
            } else {
                current[j] = max(prev[j], current[j - 1]);
            }
        }
        prev = current;
    }
    return prev[m];
}

int minInsertionsForPalindrome(string s) {
    string r = s;
    reverse(r.begin(), r.end());
    int lpsLen = lcsSpaceOptimized(s, r);
    return (int)s.size() - lpsLen;
}

int main() {
    cout << minInsertionsForPalindrome("zzazz") << endl;  // 0 (already a palindrome)
    cout << minInsertionsForPalindrome("mbadm") << endl;  // 2
    return 0;
}
\end{lstlisting}

\subsection*{Dry run}

\dryrun{Take $s = \texttt{"mbadm"}$ ($n=5$).}

The longest palindromic subsequence of \texttt{"mbadm"} is
\texttt{"mam"} (length $3$): the characters \texttt{m}, \texttt{a},
\texttt{m} appear in this order at indices $0,2,4$, and no length-$4$
palindromic subsequence exists (no other character repeats).
$\textit{minInsertions} = 5-3=2$.

\begin{complexitybox}
\textbf{Time Complexity.} $O(n^2)$, inherited from the LCS-based LPS
computation.
\textbf{Space Complexity.} $O(n)$ for the rolling-row LCS.
\end{complexitybox}

\begin{takeawaybox}
This problem closes a three-link reduction chain that is worth seeing
as a whole: Minimum Insertions for Palindrome reduces to Longest
Palindromic Subsequence, which reduces to Longest Common Subsequence.
Recognizing chains of reductions like this -- rather than treating each
problem as requiring its own fresh derivation -- is the single most
efficient way to work through a sheet of dozens of nominally distinct DP
problems.
\end{takeawaybox}

\section{Wildcard Matching}
\label{sec:dp-wildcard-matching}

\problemstatement{Given a string $s$ and a pattern $p$ containing
literal characters, \texttt{'?'} (matches any \textbf{single}
character), and \texttt{'*'} (matches any sequence of characters,
\textbf{including the empty sequence}), determine whether $p$ matches
$s$ \textbf{entirely}.}

\begin{patternbox}
This closing problem of the chapter is a deliberate capstone: it returns
to Section~\ref{sec:dp-lcs}'s match/mismatch two-pointer table, but the
``match'' condition is now governed by \textbf{pattern syntax}
(literal, \texttt{?}, or \texttt{*}) rather than simple character
equality, and \texttt{*} specifically introduces a branching choice
(match zero characters, or match one more) reminiscent of the unbounded
knapsack family's ``stay at the same index'' rule from earlier in this
book.
\end{patternbox}

\subsection*{Deriving the recurrence}

Let $f(i,j)$ be \texttt{true} if $s[0..i-1]$ (length $i$) matches
$p[0..j-1]$ (length $j$). Branch on $p[j-1]$:

\begin{itemize}
  \item If $p[j-1]$ is a literal character: it must equal $s[i-1]$
        exactly (so $i>0$ is required), and the rest must match:
        $f(i,j) = (i>0) \land (s[i{-}1]=p[j{-}1]) \land f(i{-}1,j{-}1)$.
  \item If $p[j-1] = \texttt{'?'}$: it matches any single character of
        $s$ (so $i>0$ is required, but no equality check):
        $f(i,j) = (i>0) \land f(i{-}1,j{-}1)$.
  \item If $p[j-1] = \texttt{'*'}$: it may match \textbf{zero}
        characters of $s$ (so the rest of $p$ must match all of $s$ so
        far: $f(i,j{-}1)$), \textbf{or} it may match \textbf{one more}
        character of $s$ while remaining available to match further
        characters too (so $i>0$ is required: $f(i{-}1,j)$). Either
        possibility is acceptable:
        $f(i,j) = f(i,j{-}1) \lor \big((i>0) \land f(i{-}1,j)\big)$.
\end{itemize}

Base cases: $f(0,0) = \texttt{true}$ (empty pattern matches empty
string). $f(0,j)$ for $j>0$: \texttt{true} only if $p[0..j-1]$ consists
\emph{entirely} of \texttt{'*'} characters (each can independently match
zero characters, collapsing the whole prefix to nothing).
$f(i,0)$ for $i>0$: always \texttt{false} (a non-empty string cannot be
matched by an empty pattern).

\subsection*{Approach 1: Recursion}

\begin{lstlisting}
#include <bits/stdc++.h>
using namespace std;

bool wildcardRecursive(int i, int j, string& s, string& p) {
    if (i == 0 && j == 0) return true;
    if (j == 0) return false; // i > 0 here
    if (i == 0) {
        // s is empty: only matches if remaining pattern is all '*'
        for (int k = 0; k < j; k++) if (p[k] != '*') return false;
        return true;
    }

    if (p[j - 1] == '*') {
        return wildcardRecursive(i, j - 1, s, p) ||
               wildcardRecursive(i - 1, j, s, p);
    }
    if (p[j - 1] == '?' || p[j - 1] == s[i - 1]) {
        return wildcardRecursive(i - 1, j - 1, s, p);
    }
    return false;
}

int main() {
    string s = "aa", p = "*";
    cout << (wildcardRecursive(s.size(), p.size(), s, p) ? "true" : "false") << endl;
    return 0;
}
\end{lstlisting}

\begin{complexitybox}[Approach 1 Complexity]
\textbf{Time.} $O(2^{n+m})$ in the worst case -- two branches per
\texttt{'*'} encountered.
\textbf{Space.} $O(n+m)$ recursion stack.
\end{complexitybox}

\subsection*{Approach 2: Memoization}

\begin{lstlisting}
#include <bits/stdc++.h>
using namespace std;

int wildcardMemo(int i, int j, string& s, string& p,
                  vector<vector<int>>& memo) {
    if (i == 0 && j == 0) return 1;
    if (j == 0) return 0;
    if (i == 0) {
        for (int k = 0; k < j; k++) if (p[k] != '*') return 0;
        return 1;
    }
    if (memo[i][j] != -1) return memo[i][j];

    bool result;
    if (p[j - 1] == '*') {
        result = wildcardMemo(i, j - 1, s, p, memo) ||
                  wildcardMemo(i - 1, j, s, p, memo);
    } else if (p[j - 1] == '?' || p[j - 1] == s[i - 1]) {
        result = wildcardMemo(i - 1, j - 1, s, p, memo);
    } else {
        result = false;
    }
    return memo[i][j] = result;
}

int main() {
    string s = "adceb", p = "*a*b";
    int n = s.size(), m = p.size();
    vector<vector<int>> memo(n + 1, vector<int>(m + 1, -1));
    cout << (wildcardMemo(n, m, s, p, memo) ? "true" : "false") << endl; // true
    return 0;
}
\end{lstlisting}

\begin{complexitybox}[Approach 2 Complexity]
\textbf{Time.} $O(n \cdot m)$.
\textbf{Space.} $O(n \cdot m)$ memo $+$ $O(n+m)$ recursion stack.
\end{complexitybox}

\subsection*{Approach 3: Tabulation}

\begin{lstlisting}
#include <bits/stdc++.h>
using namespace std;

bool wildcardTabulation(string& s, string& p) {
    int n = s.size(), m = p.size();
    vector<vector<bool>> dp(n + 1, vector<bool>(m + 1, false));
    dp[0][0] = true;

    for (int j = 1; j <= m; j++) {
        bool allStars = true;
        for (int k = 0; k < j; k++) if (p[k] != '*') { allStars = false; break; }
        dp[0][j] = allStars;
    }

    for (int i = 1; i <= n; i++) {
        for (int j = 1; j <= m; j++) {
            if (p[j - 1] == '*') {
                dp[i][j] = dp[i][j - 1] || dp[i - 1][j];
            } else if (p[j - 1] == '?' || p[j - 1] == s[i - 1]) {
                dp[i][j] = dp[i - 1][j - 1];
            } else {
                dp[i][j] = false;
            }
        }
    }
    return dp[n][m];
}

int main() {
    string s = "adceb", p = "*a*b";
    cout << (wildcardTabulation(s, p) ? "true" : "false") << endl; // true
    return 0;
}
\end{lstlisting}

\subsection*{Dry run of the tabulation table}

\dryrun{Take $s = \texttt{"adceb"}$, $p = \texttt{"*a*b"}$.}

Row $0$ ($s$ empty): $\textit{dp}[0][0]=\text{T}$;
$\textit{dp}[0][1]=\text{T}$ (\texttt{"*"} matches empty);
$\textit{dp}[0][2..4]=\text{F}$ (prefix \texttt{"*a"} etc.\ contain a
non-\texttt{*} character, so cannot match the empty string). Filling
forward: the first \texttt{*} absorbs \texttt{"adc"}, matching the
literal \texttt{a} against the \texttt{a} in \texttt{"adceb"}, the
second \texttt{*} absorbs nothing further, and the literal \texttt{b}
matches the trailing \texttt{b}. The final cell
$\textit{dp}[5][4]=\text{T}$, confirming the match.

\begin{complexitybox}[Approach 3 Complexity]
\textbf{Time.} $O(n \cdot m)$.
\textbf{Space.} $O(n \cdot m)$ table, no recursion stack.
\end{complexitybox}

\subsection*{Approach 4: Space Optimization}

\begin{lstlisting}
#include <bits/stdc++.h>
using namespace std;

bool wildcardSpaceOptimized(string& s, string& p) {
    int n = s.size(), m = p.size();
    vector<bool> prev(m + 1, false), current(m + 1, false);
    prev[0] = true;

    for (int j = 1; j <= m; j++) {
        bool allStars = true;
        for (int k = 0; k < j; k++) if (p[k] != '*') { allStars = false; break; }
        prev[j] = allStars;
    }

    for (int i = 1; i <= n; i++) {
        current[0] = false;
        for (int j = 1; j <= m; j++) {
            if (p[j - 1] == '*') {
                current[j] = current[j - 1] || prev[j];
            } else if (p[j - 1] == '?' || p[j - 1] == s[i - 1]) {
                current[j] = prev[j - 1];
            } else {
                current[j] = false;
            }
        }
        prev = current;
    }
    return prev[m];
}

int main() {
    string s = "adceb", p = "*a*b";
    cout << (wildcardSpaceOptimized(s, p) ? "true" : "false") << endl; // true
    return 0;
}
\end{lstlisting}

\begin{complexitybox}[Approach 4 Complexity]
\textbf{Time.} $O(n \cdot m)$.
\textbf{Space.} $O(m)$ for two rolling rows.
\end{complexitybox}

\begin{takeawaybox}
This chapter closes by combining two of its biggest ideas: the
match/mismatch two-pointer table from Section~\ref{sec:dp-lcs}, and a
branching ``match zero or match one more'' choice for \texttt{*}
strongly reminiscent of the unbounded-knapsack reuse rule from the DP on
Subsequences chapter. Recognizing that a special pattern character
introduces exactly this same kind of ``stay and try again'' branch is
the key insight that unlocks Wildcard Matching (and its sibling,
Regular Expression Matching, solved by an almost identical table) without
needing an entirely new derivation strategy.
\end{takeawaybox}

